\section{Design Decisions}\label{sec:decisions}

While building the system we had to choose between several ways of doing things. This section explains the main choices we made and why, including what we gave up.

\textbf{Keeping the interface separate from the logic.} We could have put the mark calculations directly inside the buttons of the graphical interface, which would have been quicker at first. Instead we kept the interface, the user roles, the marks classes and the file handling in separate groups of classes. This paid off because our first version ran in the console, and when we moved to the Swing interface we did not have to rewrite the marks classes. It also meant we could test the calculations from the console without opening the GUI.

\textbf{One abstract \texttt{User} class with two subclasses.} A lecturer and a Head of Department share some things, such as having a valid name, but they are allowed to do different things. We could have used one class with a variable that says which role it is and then checked it with \texttt{if} statements everywhere. We chose an abstract \texttt{User} class with \texttt{Lecturer} and \texttt{HeadOfDepartment} classes that extend it. The shared name check is written once, and only the \texttt{Lecturer} class has the methods for changing marks, so the program cannot accidentally edit marks as a Head of Department.

\textbf{Using \texttt{null} for a missing mark.} A student who has not written an assessment yet is not the same as a student who scored zero. If we stored a missing mark as $0$ or $-1$, it could be counted in the average by mistake. We store marks as \texttt{Double} objects, which can be \texttt{null}, so the program can tell that a mark has not been entered and leave that student out of the class statistics.

\textbf{Scaling each mark before applying its weight.} The simplest way to get a final mark is to multiply each raw mark by its weight. This gives wrong answers when assessments are out of different totals, for example a practical out of 50 and a test out of 100. We therefore divide each mark by its maximum mark first and then apply the weight (see Section~\ref{sec:algos}). This found a real mistake in our first version.

\textbf{Saving marks in CSV files.} We could have used a database or saved Java objects directly to a file, but both need extra tools or are hard to read. CSV files need no extra libraries, can be opened in a spreadsheet, and make importing and exporting easy. We keep one file per course. The downside is that the whole file is rewritten every time we save, and there is nothing to stop two people changing the same file at once.

\textbf{Using exceptions for bad input.} When a mark is invalid or incomplete, the marks classes throw an exception instead of quietly returning a wrong value. The \texttt{GUI} catches the exception and shows an error message, so the program keeps running. This means a wrong final mark can never be produced, but every place that asks for a result must check for missing marks first. We learned this when an early demonstration program crashed because it did not do that check (Section~\ref{sec:testing}).

\textbf{Two versions of each role method.} In Check-in 1 we agreed on a UML diagram where the \texttt{Lecturer} and \texttt{HeadOfDepartment} methods take no parameters. The GUI needs methods that take values from the screen and return text to display. Rather than change the agreed diagram, we kept the original methods and added a second version of each that takes parameters and returns text. The GUI uses the second version, so the logic classes never print anything themselves.

\textbf{A coordinator for the courses.} We added the \texttt{CourseSessionCoordinator} class to remember which course is open and which data file belongs to it. Without it, the interface would have to keep track of this itself. Now, when the user switches course, the coordinator switches the data as well.

\textbf{Trade-offs.} Plain CSV files keep the project simple and easy to inspect, but the stored data is not encrypted and there is no protection against simultaneous edits. Throwing exceptions on incomplete data makes the code a little longer, but it is safer than showing a wrong mark. We also renamed the \texttt{Module} class to \texttt{Course} so that it does not conflict with a native Java module.